%!TEX program = pdflatex
% Mathematics article -- amsart, the neutral class of the subject. Drafting
% here and converting on acceptance costs nothing, because a journal class
% changes the front matter and the page layout and leaves the body alone.
\documentclass[11pt,reqno]{amsart}

% amsart's own page is set for the AMS printed page and reads narrow on
% screen; a journal class resets it at submission.
\usepackage[letterpaper,margin=1in]{geometry}

% amsart already loads amsmath, amsthm and amsfonts.
\usepackage{amssymb,mathtools}
\usepackage{booktabs}
\usepackage{microtype}
\usepackage{bm}
\usepackage{xcolor}
% Sorts and compresses each multiple citation, so \cite{c,a,b} prints as
% [1--3] rather than in the order the keys were typed.
\usepackage{cite}
\usepackage{hyperref}
\hypersetup{colorlinks=true,
            linkcolor=blue!55!black,
            citecolor=blue!55!black,
            urlcolor=blue!55!black}
\usepackage[capitalize,nameinlink]{cleveref}

% One shared counter, so a Lemma 2.3 and a Theorem 2.4 never collide.
\theoremstyle{plain}
\newtheorem{theorem}{Theorem}[section]
\newtheorem{lemma}[theorem]{Lemma}
\newtheorem{proposition}[theorem]{Proposition}
\newtheorem{corollary}[theorem]{Corollary}
\newtheorem{conjecture}[theorem]{Conjecture}

\theoremstyle{definition}
\newtheorem{definition}[theorem]{Definition}

\theoremstyle{remark}
\newtheorem{remark}[theorem]{Remark}

\numberwithin{equation}{section}

\newcommand{\Vir}{\mathrm{Vir}}
\newcommand{\Aop}{\mathcal A}
\newcommand{\Iop}{\mathcal I}
\newcommand{\Jop}{\mathcal J}
\newcommand{\C}{\mathbb C}

\title[The identity primary at degree six]{The identity primary obstructs the
  degree-six descendant--Virasoro correspondence in the restricted
  sine--Gordon model}

\author{Zhiyuan Yao}
\address{Lanzhou Center for Theoretical Physics, Key Laboratory of Theoretical
  Physics of Gansu Province, Key Laboratory of Quantum Theory and Applications
  of MoE, Gansu Provincial Research Center for Basic Disciplines of Quantum
  Physics, Lanzhou University, Lanzhou 730000, China}
\email{yaozy@lzu.edu.cn}

\date{\today}

\subjclass[2020]{Primary 81T40; Secondary 17B68, 81R12}
\keywords{restricted sine--Gordon model, Virasoro module, descendants,
  null vectors, local integrals of motion, vacuum module}

\begin{document}

\begin{abstract}
The local fields of the restricted sine--Gordon model carry two descriptions.
One acts on towers of form-factor polynomials by the odd local integrals of
motion and by transverse generators; the other organizes descendants into
Virasoro modules. These are conjectured to agree degree by degree, and at
degree six the comparison sets eleven descendants of the primary
$\phi_{2m}$ against the eleven Virasoro descendants indexed by the partitions
of six, for every $m\ge0$ and every allowed coupling. We evaluate the
endpoint $m=0$. There the primary is the identity operator and the local
integrals act by commutators, so each of them annihilates it; eight of the
eleven displayed descendants are products whose rightmost factor is such an
integral, and therefore vanish. The remaining span has dimension at most
three. The universal vacuum Verma module has degree-six dimension eleven, and
its quotient by the submodule generated by the singular vector $L_{-1}v$ has
dimension four. At the coupling $\xi=\sqrt2$, which satisfies the
non-root-of-unity condition, the asserted equality consequently fails under
either target. The failure is sharp and lands on the evidence advanced for the
identification: that evidence is a count, matching the operators which vanish
against the null vectors of the module, and at this degree eight operators
vanish where seven null vectors are removed. A formulation that survives must
exclude the identity sector, prescribe a limiting normalization there, or name
a target of dimension at most three; the cases $m\ge1$ are untouched by this
argument.
\end{abstract}

\maketitle

\section{Introduction}
\label{sec:introduction}

Two-dimensional conformal field theory organizes the local fields of a model
into Virasoro modules, with descendants of a primary field labelled by
products of negative Virasoro modes \cite{Belavin_1984ic}. A massive
integrable perturbation of such a theory retains an infinite family of local
integrals of motion \cite{Zamolodchikov_1987hi,Zamolodchikov_1989if}, whose
integrable structure in the conformal limit was developed by Bazhanov,
Lukyanov and Zamolodchikov \cite{Bazhanov_1996is,Bazhanov_1997is,
Bazhanov_1999is}, and its local
fields are instead produced by the form-factor bootstrap, as towers of
polynomials in the particle rapidities \cite{Smirnov_Book}. The restricted
sine--Gordon model, obtained by restricting the sine--Gordon solitons at a
root of unity so that the ultraviolet limit is a minimal model
\cite{Leclair_1989rs,Smirnov_1990ro,Bernard_1990rq}, is the setting where the
two descriptions have been compared most closely.

Smirnov proposed that the comparison is an identification, and related the
null relations on the form-factor side to deformed Abelian differentials
\cite{Smirnov_1996ot,Smirnov_1998so}. Concretely, one fixes $\xi>0$ with
$q=\exp(2\pi i/\xi)$ not a root of unity, takes the primary $\phi_{2m}$ of
weight $\Delta_{2m}=-m+m(m+1)\xi$ at central charge
$c=1-6/(\xi(\xi+1))$, and generates descendants with two families of
operators: the odd local integrals of motion $I_{2k-1}$, acting on a local
operator $\mathcal O$ by the commutator $\Iop_{2k-1}\mathcal O=[I_{2k-1},
\mathcal O]$ and carrying degree $2k-1$, and transverse generators
$\Jop_{2k}$ of degree $2k$, acting on a form-factor polynomial by
multiplication. The assertion is that the span of such descendants of a given
degree coincides with the corresponding graded piece of the Virasoro module
generated by $\phi_{2m}$. Babelon, Bernard, and Smirnov developed the
matching null-vector theory, identified the vacuum primary with the identity
operator, and derived the character
$(1-p^{2m+1})\prod_{j\ge1}(1-p^{j})^{-1}$ of the resulting quotient
\cite{Babelon_1997nv}.

The evidence Smirnov offers for the identification is a count. In the
discussion following his Conjecture he observes first that the two graded
spaces have the same character, and then states as the crucial check that the
local operators which vanish are exactly as numerous as the null vectors of
the Verma module, crediting that count to the null-vector paper
\cite{Smirnov_1998so,Babelon_1997nv}.
That check is what the present note evaluates at one point, and it is where
the failure appears: at degree six in the identity sector the two numbers
differ by one.

Later fermionic constructions made the correspondence explicit while keeping a
sharp boundary in degree. The hidden Grassmann structure of the XXZ chain and
its conformal and sine--Gordon limits supplied a fermionic basis
\cite{Boos_2007hg,Boos_2008hg,Jimbo_2009hg,Boos_2010hg,Jimbo_2010hg}, and the
identification of that basis with the Virasoro one was carried through level
six, then level eight, and then level ten, in every case working modulo the
action of the local integrals of motion
\cite{Boos_2010hg,Boos_2011fb,Negro_2013rr}. Jimbo, Miwa, and Smirnov compared
fermionic and Virasoro descendants in the same modulo-integrals sense
\cite{Jimbo_2011fs}. Removing that quotient is what makes the problem hard.
Boos and Smirnov achieved it with nontrivial integral action through level
four; at level six they report that the eleven descendants yield only nine
linearly independent polynomials in their expectation-value scheme, so
boundary data from ground states alone do not determine the coefficients, and
they defer the case \cite{Boos_2018nr}.

The degree-six statement left open there is the subject of this note. Written
out, it asserts for every $m\ge0$ and every allowed coupling that the span
of
\begin{equation}
\begin{gathered}
\Jop_6\phi_{2m},\quad
\Iop_5\Iop_1\phi_{2m},\quad
\Jop_4\Jop_2\phi_{2m},\quad
\Jop_4\Iop_1^2\phi_{2m},\quad
\Iop_3^2\phi_{2m},\quad
\Iop_3\Jop_2\Iop_1\phi_{2m},\\
\Iop_3\Iop_1^3\phi_{2m},\quad
\Jop_2^3\phi_{2m},\quad
\Jop_2^2\Iop_1^2\phi_{2m},\quad
\Jop_2\Iop_1^4\phi_{2m},\quad
\Iop_1^6\phi_{2m}
\end{gathered}
\label{eq:operator-list}
\end{equation}
equals the degree-six subspace of the Virasoro module generated by
$\phi_{2m}$, spanned by the eleven monomials $L_{-6}\phi_{2m}$,
$L_{-5}L_{-1}\phi_{2m},\ldots,L_{-1}^6\phi_{2m}$ indexed by the partitions of
six.

Before a transition matrix between two eleven-element lists is sought, the
quantifier deserves attention: $m\ge0$ includes $m=0$, and there the primary
is the identity operator. Our observation is that this endpoint already
settles the question, in the negative and by an exact count.

\begin{theorem}
\label{thm:main-intro}
Take $\xi=\sqrt2$, so that $q=\exp(2\pi i/\xi)$ is not a root of unity, and
take $m=0$, so that $\phi_0=1$. Let $\Aop_0[6]$ be the span of the eleven
operators in \eqref{eq:operator-list} at $m=0$. Then
\[
  \dim\Aop_0[6]\le3,
\]
whereas the degree-six subspace of the universal Verma module $M(c,0)$ has
dimension $11$ and the degree-six subspace of its quotient by the submodule
generated by the singular vector $L_{-1}v$ has dimension $4$. Consequently
$\Aop_0[6]$ equals neither.
\end{theorem}

The mechanism is a single identity. The action of an integral of motion is a
commutator, and a commutator with the identity operator vanishes, so
$\Iop_{2k-1}\phi_0=0$ for every $k$. Operator products in
\eqref{eq:operator-list} act on the primary from the right, so any monomial
whose rightmost factor is an $\Iop$ is annihilated. Seven of the eleven end
in $\Iop_1$ and an eighth, $\Iop_3^2\phi_0$, acts first by $\Iop_3$; only the
three pure-$\Jop$ monomials can survive. Three is below four, and four is
already the smaller of the two conventional target dimensions, so the
mismatch does not turn on which Virasoro target is meant.

This is a statement about the formulation, not a refutation of the programme
it belongs to. What it shows is that the identity sector is a genuine
constraint on how the correspondence may be posed: a target must be specified
together with its behaviour at $m=0$, since the operator side degenerates
there for reasons that have nothing to do with the coupling or with degree
six in particular. It says nothing about $m\ge1$, where the displayed
descendants need not degenerate and the eleven-against-eleven comparison
retains its content.

\Cref{sec:formulation} fixes the conventions and the two candidate targets.
\Cref{sec:vacuum} proves the two facts about the identity sector.
\Cref{sec:rank} proves \cref{thm:main-intro}. \Cref{sec:boundary} delimits
what the theorem decides and states the formulations that survive it.

\section{The degree-six comparison and its two Virasoro targets}
\label{sec:formulation}

Fix $\xi>0$ and set
\begin{equation}
  q=\exp(2\pi i/\xi),
  \qquad q\ \text{not a root of unity}.
  \label{eq:q}
\end{equation}
For $m\ge0$, let $\phi_{2m}$ be the restricted sine--Gordon primary of weight
\begin{equation}
  \Delta_{2m}=-m+m(m+1)\xi,
  \qquad
  c=1-\frac{6}{\xi(\xi+1)}.
  \label{eq:weights}
\end{equation}
The two families of descendant-generating operators carry degrees
$\deg\Iop_{2k-1}=2k-1$ and $\deg\Jop_{2k}=2k$. The first acts by the
commutator with the local integral of motion of spin $2k-1$, a definition
stated in this form by Smirnov and used unchanged since
\cite{Smirnov_1998so,Jimbo_2011fs},
\begin{equation}
  \Iop_{2k-1}\mathcal O=[I_{2k-1},\mathcal O],
  \label{eq:commutator}
\end{equation}
and the second acts on the $2n$-particle form-factor polynomial
$L_n[\mathcal O](A\mid B)$ by
\begin{equation}
 L_n[\Jop_{2k}\mathcal O](A\mid B)
 =\Bigl(\sum_{i=1}^{n}A_i^{2k}-\tfrac12\sum_{j=1}^{2n}B_j^{2k}\Bigr)
 L_n[\mathcal O](A\mid B).
 \label{eq:transverse}
\end{equation}
Write $\Aop_m[d]$ for the span, inside the space of local operators, of the
degree-$d$ descendants of $\phi_{2m}$ generated this way. Only
\eqref{eq:commutator} enters the argument below; \eqref{eq:transverse} is
recorded because it fixes which of the eleven monomials are pure-$\Jop$, and
no property of $\Jop_{2k}$ beyond its degree is used.

\begin{remark}
\label{rem:null}
The polynomial multiplier representing $\Iop_{2k-1}$ before the form-factor
pairing need not vanish identically, and \eqref{eq:commutator} does not say
that it does. The two are consistent because the polynomial presentation has a
null ideal: a tower all of whose form factors vanish represents the zero local
operator \cite{Babelon_1997nv}. The identity \eqref{eq:commutator} is read in
the space of local operators, which is where the asserted equality of spans
lives.
\end{remark}

The degree-six operator span $\Aop_m[6]$ is the span of the eleven monomials
listed in \eqref{eq:operator-list}. On the Virasoro side the comparison list
is indexed by the eleven partitions of six,
\begin{equation}
\begin{gathered}
L_{-6}\phi_{2m},\quad L_{-5}L_{-1}\phi_{2m},\quad
L_{-4}L_{-2}\phi_{2m},\quad L_{-4}L_{-1}^2\phi_{2m},\quad
L_{-3}^2\phi_{2m},\quad L_{-3}L_{-2}L_{-1}\phi_{2m},\\
L_{-3}L_{-1}^3\phi_{2m},\quad L_{-2}^3\phi_{2m},\quad
L_{-2}^2L_{-1}^2\phi_{2m},\quad L_{-2}L_{-1}^4\phi_{2m},\quad
L_{-1}^6\phi_{2m},
\end{gathered}
\label{eq:vir-list}
\end{equation}
where the modes satisfy
\begin{equation}
  [L_r,L_s]=(r-s)L_{r+s}+\frac{c}{12}(r^3-r)\delta_{r+s,0}.
  \label{eq:virasoro}
\end{equation}

The phrase ``the Virasoro module generated by $\phi_{2m}$'' admits two
readings, and we keep both. Let $M(c,h)$ be the universal highest-weight
module induced from a vector $v$ with $L_0v=hv$ and $L_nv=0$ for $n>0$. By
the Poincar\'e--Birkhoff--Witt theorem the ordered monomials
\begin{equation}
  L_{-a_1}L_{-a_2}\cdots L_{-a_k}v,
  \qquad a_1\ge a_2\ge\cdots\ge a_k\ge1,
  \label{eq:pbw}
\end{equation}
form a basis of $M(c,h)$, so a singular vector is a nonzero element of the
universal module and becomes zero only in a quotient
\cite{Feigin_1984vm,Kac_cf}. At $h=0$ the vector $L_{-1}v$ is singular:
$L_1L_{-1}v=2L_0v=0$ by \eqref{eq:virasoro}, and $L_nL_{-1}v=(n+1)L_{n-1}v=0$
for $n>1$. The second reading therefore removes it, giving the
translation-invariant vacuum module
\begin{equation}
  V^{\mathrm{vac}}(c)=M(c,0)\big/U(\Vir_-)L_{-1}v,
  \label{eq:vacuum-quotient}
\end{equation}
which is the standard vacuum module of a chiral algebra
\cite{Rocha-Caridi_1985vv}.

Both readings have support in the primary literature, which is why the
argument below is run against both. Smirnov states the correspondence using
the words ``Verma module'' while separately distinguishing the reducible
Verma module from the irreducible representation obtained by factoring out its
singular submodule \cite{Smirnov_1998so}. Babelon, Bernard, and Smirnov work
with the physical quotient, identifying the vacuum primary with the identity
operator and computing the character
$(1-p^{2m+1})\prod_{j\ge1}(1-p^{j})^{-1}$, whose $m=0$ specialization
$(1-p)\prod_{j\ge1}(1-p^{j})^{-1}$ is the character of
\eqref{eq:vacuum-quotient} \cite{Babelon_1997nv}.

\section{The identity sector}
\label{sec:vacuum}

\begin{lemma}[Admissible coupling]
\label{lem:allowed}
The value $\xi=\sqrt2$ satisfies the condition in \eqref{eq:q}: with
$q=\exp(2\pi i/\sqrt2)$, there is no integer $N>0$ with $q^N=1$.
\end{lemma}

\begin{proof}
Suppose $q^N=1$ for an integer $N>0$. Then $2\pi iN/\sqrt2\in2\pi i\mathbb Z$,
so $N/\sqrt2=r$ for some $r\in\mathbb Z$, and $r\ne0$ because $N>0$. Hence
$\sqrt2=N/r\in\mathbb Q$, contradicting the irrationality of $\sqrt2$.
\end{proof}

\begin{lemma}[Vacuum commutator]
\label{lem:vacuum}
At $m=0$ the primary is the identity operator, $\phi_0=1$, and
\begin{equation}
  \Iop_{2k-1}\phi_0=0
  \qquad\text{for every }k\ge1.
  \label{eq:vacuum-zero}
\end{equation}
\end{lemma}

\begin{proof}
The identification $\phi_0=1$ is the vacuum convention of the null-vector
construction, where the primary of weight $\Delta_0=0$ is the identity field
\cite{Babelon_1997nv}. The conclusion is then immediate from
\eqref{eq:commutator}: the identity operator commutes with every operator, so
$\Iop_{2k-1}\phi_0=[I_{2k-1},1]=0$.
\end{proof}

The weight is consistent: \eqref{eq:weights} gives $\Delta_0=0$ at $m=0$ for
every $\xi$, so the specialization does not depend on \cref{lem:allowed}, and
\cref{lem:allowed} is needed only to place the pair $(\xi,m)=(\sqrt2,0)$
inside the range the assertion quantifies over.

\Cref{lem:vacuum} is not a reading imposed from outside. The same vanishing is
recorded in the null-vector construction itself, by a different route: writing
the vacuum primary as $\Phi_0=1$, Babelon, Bernard, and Smirnov list the basic
null vector of its module as $\Gamma_0=L_{-1}\Phi_0$, and identify the simplest
null vector of $\Phi_0$ with a multiple of $I_1\Phi_0$ \cite{Babelon_1997nv}. A
null vector is an operator all of whose form factors vanish, hence the zero
local operator, so $\Iop_1\phi_0=0$ is already their statement. What
\eqref{eq:commutator} adds is that the same holds for every $k$, at once and
without computing a form factor.

\begin{remark}
\label{rem:degree-one}
The same identity already separates the two sides at degree one, where the
only operator of degree one is $\Iop_1$, so $\Aop_0[1]=0$, while
\eqref{eq:pbw} gives $M(c,0)[1]=\C L_{-1}v$ of dimension one. That mismatch
is removed by passing to \eqref{eq:vacuum-quotient}, whose degree-one
subspace is zero. The point of the next section is that at degree six no such
repair is available.
\end{remark}

\section{The rank obstruction at degree six}
\label{sec:rank}

\begin{proposition}[Operator-side rank]
\label{prop:rank}
At $m=0$, eight of the eleven monomials in \eqref{eq:operator-list} vanish,
and
\begin{equation}
  \dim\Aop_0[6]\le3.
  \label{eq:rank-three}
\end{equation}
\end{proposition}

\begin{proof}
An operator product in \eqref{eq:operator-list} acts on $\phi_0$ from the
right, so if its rightmost factor is some $\Iop_{2k-1}$ then the whole
product annihilates $\phi_0$ by \cref{lem:vacuum}, whatever the remaining
factors are. \Cref{tab:classification} records the rightmost factor of each
of the eleven monomials. Seven of them end in $\Iop_1$, namely
$\Iop_5\Iop_1$, $\Jop_4\Iop_1^2$, $\Iop_3\Jop_2\Iop_1$, $\Iop_3\Iop_1^3$,
$\Jop_2^2\Iop_1^2$, $\Jop_2\Iop_1^4$ and $\Iop_1^6$; an eighth, $\Iop_3^2$,
ends in $\Iop_3$. All eight therefore give the zero operator. The remaining
three are $\Jop_6\phi_0$, $\Jop_4\Jop_2\phi_0$ and $\Jop_2^3\phi_0$, and a
span of three vectors has dimension at most three. No claim about the
independence of those three is needed, and none is made.
\end{proof}

\begin{table}[t]
\caption{The eleven degree-six operator monomials of
\eqref{eq:operator-list} at $m=0$, classified by the rightmost factor
through which each acts on $\phi_0$. Every monomial whose rightmost factor is
an integral of motion is annihilated by \cref{lem:vacuum}; the three
pure-$\Jop$ monomials are the only ones not forced to vanish.}
\label{tab:classification}
\begin{tabular}{llc}
\toprule
Monomial & Rightmost factor & Value at $m=0$ \\
\midrule
$\Jop_6\phi_0$                 & $\Jop_6$  & not forced to vanish \\
$\Jop_4\Jop_2\phi_0$           & $\Jop_2$  & not forced to vanish \\
$\Jop_2^3\phi_0$               & $\Jop_2$  & not forced to vanish \\
\midrule
$\Iop_5\Iop_1\phi_0$           & $\Iop_1$  & $0$ \\
$\Jop_4\Iop_1^2\phi_0$         & $\Iop_1$  & $0$ \\
$\Iop_3\Jop_2\Iop_1\phi_0$     & $\Iop_1$  & $0$ \\
$\Iop_3\Iop_1^3\phi_0$         & $\Iop_1$  & $0$ \\
$\Jop_2^2\Iop_1^2\phi_0$       & $\Iop_1$  & $0$ \\
$\Jop_2\Iop_1^4\phi_0$         & $\Iop_1$  & $0$ \\
$\Iop_1^6\phi_0$               & $\Iop_1$  & $0$ \\
$\Iop_3^2\phi_0$               & $\Iop_3$  & $0$ \\
\bottomrule
\end{tabular}
\end{table}

\begin{proposition}[The two target dimensions]
\label{prop:dimensions}
At degree six,
\begin{equation}
  \dim M(c,0)[6]=11,
  \qquad
  \dim V^{\mathrm{vac}}(c)[6]=4.
  \label{eq:dimensions}
\end{equation}
\end{proposition}

\begin{proof}
By \eqref{eq:pbw} the monomials $L_{-a_1}\cdots L_{-a_k}v$ with
$a_1\ge\cdots\ge a_k\ge1$ and $\sum_ia_i=6$ form a basis of $M(c,0)[6]$, so
its dimension is the number of partitions of six, $p(6)=11$.

For the quotient, the degree-six subspace of $U(\Vir_-)L_{-1}v$ is spanned by
the vectors $XL_{-1}v$ with $X\in U(\Vir_-)$ of degree five, and $U(\Vir_-)$
in degree five is spanned by ordered monomials
$L_{-a_1}\cdots L_{-a_k}$ with $a_1\ge\cdots\ge a_k\ge1$ and
$\sum_ia_i=5$. For such a monomial the product
$L_{-a_1}\cdots L_{-a_k}L_{-1}v$ is again ordered, because $a_k\ge1$; it is
therefore already one of the basis vectors \eqref{eq:pbw}, namely the one
indexed by the partition of six obtained by appending a part equal to $1$.
That map from partitions of five to partitions of six is a bijection onto
those partitions of six having at least one part equal to $1$, so the
spanning set consists of $p(5)=7$ distinct basis vectors of $M(c,0)[6]$ and
is in particular linearly independent. Hence
$\dim\bigl(U(\Vir_-)L_{-1}v\bigr)[6]=7$ exactly, and
\[
  \dim V^{\mathrm{vac}}(c)[6]=p(6)-p(5)=11-7=4.
\]
The surviving quotient basis is indexed by the four partitions of six with no
part equal to one: $6$, $4+2$, $3+3$, and $2+2+2$. The same number is the
coefficient of $p^6$ in the character
$(1-p)\prod_{j\ge1}(1-p^{j})^{-1}$ of \cite{Babelon_1997nv} at $m=0$.
\end{proof}

\begin{proof}[Proof of \cref{thm:main-intro}]
By \cref{lem:allowed} the coupling $\xi=\sqrt2$ satisfies \eqref{eq:q}, so
the pair $(\xi,m)=(\sqrt2,0)$ is one the assertion quantifies over.
\Cref{prop:rank} bounds $\dim\Aop_0[6]$ by $3$, and \cref{prop:dimensions}
gives the two target dimensions $11$ and $4$. Since $3<4<11$, the space
$\Aop_0[6]$ has dimension different from both, so it equals neither.
\end{proof}

\begin{corollary}[The counting check fails by one]
\label{cor:count}
At $m=0$ and degree six, the number of displayed monomials forced to vanish is
eight, while the number of null vectors removed from the Verma module at that
degree is $p(5)=7$. The two counts that Smirnov's crucial check requires to
agree therefore differ by one.
\end{corollary}

\begin{proof}
\Cref{prop:rank} exhibits eight monomials of \eqref{eq:operator-list}
annihilated by \cref{lem:vacuum}. The proof of \cref{prop:dimensions} shows
that the degree-six part of $U(\Vir_-)L_{-1}v$, the submodule of null vectors
at $m=0$, has dimension exactly $p(5)=7$. If the span of the eleven monomials
were to fill the quotient \eqref{eq:vacuum-quotient}, whose degree-six
dimension is $11-7=4$, then at most seven of them could vanish. Eight do.
\end{proof}

Stated this way the obstruction is not a mismatch with a convention but a
failure of the evidence advanced for the correspondence. Smirnov supports the
conjecture by observing that the operators which vanish are as numerous as the
null vectors \cite{Smirnov_1998so}. At $m=0$ they are more
numerous, and the surplus is exactly the one dimension by which the operator
span falls short of the vacuum quotient. The discrepancy is small, which is
part of why it survives a character comparison: characters agree degree by
degree only when the vanishing is counted correctly, and at this point it is
not.

The theorem is stronger than exhibiting a linear dependence among the eleven
operator candidates. A dependence would reduce the operator side to some
dimension below eleven and leave open whether the physical quotient absorbs
the difference. Here the shortfall survives the quotient: even after the
vacuum singular vector and all of its descendants are removed, the target
retains four dimensions against an operator span of at most three.

\begin{remark}[On computation]
\label{rem:verification}
No step above is computer assisted. The classification in
\cref{tab:classification} is an inspection of eleven written monomials, and
the two dimension counts are the partition numbers $p(6)=11$ and $p(5)=7$;
both can be checked by hand in a few minutes, and both are stated in the text
rather than delegated. We did check the transcription of the eleven monomials,
the eight forced zeros, and the two partition counts with an independent
short program, but it certifies a transcription rather than an argument, and
the proof does not rest on it.
\end{remark}

\section{Scope, and the formulations that survive}
\label{sec:boundary}

\Cref{thm:main-intro} rests on three readings of the assertion, and it is
worth being explicit that they are readings. The quantifier $m\ge0$ is taken
to include $m=0$; the primary there is the identity operator without a
limiting renormalization; and the operator $\Iop_{2k-1}$ acts by the
commutator \eqref{eq:commutator}. Each is what the statement says, and the
three are mutually consistent, which is what makes the finite argument exact.
Together they are also the whole content of the hypothesis: the argument uses
no property of $\Jop_{2k}$, no value of the central charge, and no feature of
degree six beyond the list \eqref{eq:operator-list} itself.

What the theorem does not establish is that the authors of the correspondence
intended this endpoint. The primary sources mix the two levels of description
deliberately. Smirnov phrases the proposal with ``Verma module'' while
elsewhere separating that module from its irreducible quotient, and discusses
vanishing local operators in the same passage \cite{Smirnov_1998so}; Babelon,
Bernard, and Smirnov work throughout with the physical quotient and its
character \cite{Babelon_1997nv}. Their printed coupling, weight, deformation,
and integral normalizations do not admit a verified term-by-term transport
into every later convention, so we have not attempted one. The argument above
needs none, since it uses only the commutator identity and the displayed
list.

Nor does the theorem constrain $m\ge1$. There $\phi_{2m}$ is not the identity,
$\Iop_{2k-1}\phi_{2m}$ need not vanish, and the eleven-against-eleven
comparison keeps its content; the rank shortfall proved here is a property of
the identity sector alone.

Several mathematically distinct formulations survive, and they are genuinely
different problems rather than restatements of one another:
\begin{enumerate}
  \item restrict the assertion to $m\ge1$, and establish the degree-six
    transition matrix for generic non-root-of-unity coupling;
  \item prescribe a renormalized $m\to0$ limit in which the displayed
    descendants are rescaled before the primary is specialized, so that the
    identity sector is approached rather than evaluated;
  \item define the null ideal of the form-factor presentation intrinsically,
    and prove that the resulting quotient agrees with the appropriate
    Virasoro singular quotient; or
  \item replace equality by an equality modulo the local integrals of motion,
    keeping explicit track of the discarded directions.
\end{enumerate}
The fourth is the sense in which the correspondence is already established
through level six and beyond \cite{Boos_2010hg,Boos_2011fb,Negro_2013rr,
Jimbo_2011fs}, and it is consistent with \cref{thm:main-intro} for a concrete
reason: quotienting by the integrals of motion removes precisely the
directions whose absence the theorem counts. Exact
expectation values without that quotient are known only at low level: at
level two for the Virasoro descendant $L_{-2}\bar L_{-2}\Phi$
\cite{Fateev_1999ev}, and at level four once the integrals of motion act
nontrivially \cite{Boos_2018nr}; see also the study of one-point functions of
descendants in the sine-Gordon model \cite{Jimbo_2010oo}. Degree six is the next case, and it is where
the present question begins.

One repair is not available. An unspecified ``physical quotient'' cannot
rescue the literal statement by stipulation, because the theorem rules out
both quotients that the sources actually define. Any smaller target must be
specified independently, and its degree-six dimension and the operator map
into it computed; a target of dimension at most three is not excluded by this
argument, but neither is it supplied by it.

\section{Conclusion}
\label{sec:conclusion}

The identity sector is a sharp test of any proposed correspondence between
integral-of-motion descendants and Virasoro descendants, because the operator
side degenerates there for a reason internal to the definition: the action is
a commutator, and the identity operator commutes with everything. Applied to
the degree-six list, the test is decisive. Eight of the eleven descendants
vanish, the span drops to dimension at most three, and the two conventional
Virasoro targets retain dimensions eleven and four. At $\xi=\sqrt2$ and
$m=0$, an allowed pair, the asserted equality fails.

What makes the failure informative rather than merely negative is where it
lands. The correspondence is supported by a count of vanishing operators
against null vectors, and \cref{cor:count} shows that count to be wrong by one
at this degree and this primary. A single dimension is small enough to survive
a character comparison and large enough to defeat an equality of spans, which
is the gap between the two kinds of evidence.

The consequence is a constraint on how the correspondence is posed rather
than on whether it holds. A degree-six transition matrix can be claimed only
once the target module and the treatment of the endpoint $m=0$ are stated
together, since the two interact: the quotient that repairs degree one, as in
\cref{rem:degree-one}, does not repair degree six. The natural next step is
the first of the surviving formulations above --- the assertion restricted to
$m\ge1$, where the eleven-against-eleven comparison is not degenerate and
where the reported linear dependence among the level-six expectation-value
polynomials \cite{Boos_2018nr} suggests the answer is not simply affirmative
either. Settling the intended convention at $m=0$ would determine whether
this note records a counterexample or a needed hypothesis.

\section*{Data and code availability}

The proof is analytic and finite, and no data underlie it. The short program
described in \cref{rem:verification}, which independently re-checks the
transcription of the eleven monomials, the eight forced zeros, and the
partition counts $p(6)=11$ and $p(5)=7$, is available from the author on
request. It has no bearing on the validity of the argument, which is given in
full in \cref{sec:rank}.

\section*{Acknowledgments}

Z.\ Y.\ acknowledges support from the Strategic Priority Research Program of
the Chinese Academy of Sciences under Grant No.~XDB1680102. Z.\ Y.\ also
acknowledges support by National Natural Science Foundation of China (Grant
No.~12247101). We also acknowledge Gewu Intelligence Lab for providing the
collaborative research environment in which this project was developed.

OpenAI GPT-5.6 Sol was used substantively in this work: in solving the
problem, including the formulation of arguments and the derivation of
intermediate steps, and in drafting and revising the manuscript. The author
specified the problem, guided the approach to be taken, checked every argument
and result it produced, and takes full responsibility for the content of this
manuscript. In particular, the classification in \cref{tab:classification},
the partition counts in \cref{prop:dimensions}, and every statement
attributed to a cited work were checked by hand against the definitions and
against the cited sources.

\bibliographystyle{unsrt}
\bibliography{references,unverified}

\end{document}
